% SGA 3, Expose VI_B, Loop 1 English reconstruction.
% Scope: Proposition 2.5, Corollary 2.6, Proposition 2.7, and
% Lemma 2.7.1 in full; Corollary 2.8 is excluded.
% Authority: Polo--Gille Exp6B-13oct24.pdf, component-local pages 11--14
% (combined-reader pages 352--355; original-margin printed pages 333--336).
% OCR and the Reinhold English Markdown are comparison witnesses only.

\medskip
\phantomsection\label{prop:sga3-VIB-2.5}%
\noindent\textbf{命题 2.5.---}%
\footnote{\label{note:sga3-VIB-20}编者注: 展开了陈述及证明,
并在 \textup{\SGARefLink{concl:sga3-VIB-2.5-i}{(i)}} 中将 \(H\)
的假设从``有限呈示''减弱为``有限型''.}
\emph{设 \(u:G\to H\) 为 \(S\)-群态射.}

\begin{enumerate}
\item[\textup{(i)}]\phantomsection\label{concl:sga3-VIB-2.5-i}
  \emph{假设 \(G\) (分别地, \(H\)) 在其单位截面各点处
  在 \(S\) 上为有限呈示且平坦 (分别地, 有限型). 则由使 \(u_s\)
  平坦 (分别地, 平滑或平展) 的点 \(s\in S\) 组成的集合}
  \[
    U_{\mathrm{plat}}
      \supset U_{\mathrm{lisse}}
      \supset U_{\mathrm{\acute{e}t.}},
  \]
  \emph{在 \(S\) 中开.}

  \emph{若此外 \(G\) (分别地, \(H\)) 在 \(S\) 上局部有限呈示且平坦
  (分别地, 局部有限型), 则 \(G\) 中使 \(u\) 平坦
  (分别地, 平滑或平展) 的点集 \(V_{\mathrm{plat}}\)
  (分别地, \(V_{\mathrm{lisse}}\) 或 \(V_{\mathrm{\acute{e}t.}}\))
  是 \(U_{\mathrm{plat}}\) (分别地, \(U_{\mathrm{lisse}}\)
  或 \(U_{\mathrm{\acute{e}t.}}\)) 在 \(\pi_G\) 下的逆像.}

\item[\textup{(ii)}]\phantomsection\label{concl:sga3-VIB-2.5-ii}
  \emph{假设对每个 \(s\in S\), \(G_s\) 在 \(\kappa(s)\) 上局部有限型
  (若 \(G\) 在单位截面各点处在 \(S\) 上为有限型,
  此条件成立 \textup{(\SGARefLink{lem:sga3-VIB-1.3.1.1}{1.3.1.1})}),
  且 \(u\) 在 \(G\) 的单位截面各点处局部有限型
  (分别地, 局部有限呈示). 则由使 \(u_s\) 局部拟有限
  (分别地, 非分歧) 的 \(s\in S\) 组成的集合}
  \[
    U_{\mathrm{l.q.f.}}\supset U_{\mathrm{n.r.}},
  \]
  \emph{在 \(S\) 中开.}

  \emph{若此外 \(u\) 局部有限型 (分别地, 局部有限呈示),
  则 \(G\) 中使 \(u\) 拟有限 (分别地, 非分歧) 的点集
  \(V_{\mathrm{l.q.f.}}\) (分别地, \(V_{\mathrm{n.r.}}\))
  是 \(U_{\mathrm{l.q.f.}}\) (分别地, \(U_{\mathrm{n.r.}}\))
  在 \(\pi_G\) 下的逆像.}
\end{enumerate}

证明 \textup{\SGARefLink{concl:sga3-VIB-2.5-i}{(i)}}.
先注意 (\SGARefLink{lem:sga3-VIB-1.3.1.1}{1.3.1.1}),
对每个 \(s\in S\), \(G_s\) 在 \(\kappa(s)\) 上局部有限型.
设 \(Y\) 为 \(H\) 中使 \(\pi_H\) 为有限型的点组成的开子集,
\(X_1\) 为 \(u^{-1}(Y)\) 中使 \(\pi_G\) 为有限呈示的点组成的开子集.
由 EGA~IV\(_3\), 11.3.1,
\(X_1\) 中使 \(\pi_G\) 为有限呈示且平坦的点集 \(X\)
在 \(X_1\) 中开, 因而在 \(G\) 中开.

设
\[
  u_X:X\longrightarrow Y
\]
为 \(u\) 诱导的态射. 由于 \(Y\) (分别地, \(X\))
包含 \(H\) (分别地, \(G\)) 的单位截面,
\(\pi_G\) 在 \(X\) 上的限制 \(\pi_X\) 满射,
且 \(\epsilon_G\) 诱导的态射
\[
  S\longrightarrow X
\]
是 \(X\) 的 \(S\)-截面.

考虑 \(X\) 中使 \(u_X\) 平坦 (分别地, 平滑或平展)
的点组成的集合
\[
  W_{\mathrm{plat}}
    \supset W_{\mathrm{lisse}}
    \supset W_{\mathrm{\acute{e}t.}},
\]
显然, \(W_{\mathrm{lisse}}\) 和 \(W_{\mathrm{\acute{e}t.}}\)
在 \(X\) 中开; 参见编者\hyperref[note:sga3-VIB-6]{注~6}.
另一方面, 由于 \(\pi_Y\) 局部有限型, 且
\[
  \pi_X=\pi_Y\circ u_X
\]
局部有限呈示, EGA~IV\(_1\), 1.4.3(v) 表明 \(u_X\) 局部有限呈示.
因此, \(u_X\) 的平坦点集 \(W_X\) 在 \(X\) 中开
(EGA~IV\(_3\), 11.3.1).

设 \(x\in X\), 令 \(s=\pi_X(x)\).
由 EGA~IV\(_2\), 11.3.10,%
\begingroup
\renewcommand{\thefootnote}{\fnsymbol{footnote}}%
\footnote[1]{\label{note:sga3-VIB-2.5-EGA-volume}译者注:
保留了底本的引文. 第 11 节位于 EGA~IV\(_3\), 而非 EGA~IV\(_2\);
所引结果为 EGA~IV\(_3\), 11.3.10.}%
\endgroup{}
结合 EGA~IV\(_4\), 17.5.1 (分别地, 17.6.1),
\(x\) 属于 \(W_{\mathrm{plat}}\)
(分别地, \(W_{\mathrm{lisse}}\) 或 \(W_{\mathrm{\acute{e}t.}}\))
当且仅当 \(u_s\) 在点 \(x\) 处平坦 (分别地, 平滑或平展),
或由 \hyperref[prop:sga3-VIB-1.3]{1.3}, 等价地,
当且仅当 \(u_s\) 平坦 (分别地, 平滑或平展).

因此, 对于
\[
  \text{``\(\flat=\) 平坦、平滑或平展,''}
\]
有
\[
  U_\flat=\epsilon_G^{-1}(W_\flat),
  \qquad
  W_\flat=\pi_X^{-1}(U_\flat).
\]
由于 \(W_\flat\) 在 \(X\) 中开, 因而在 \(G\) 中开,
第一个等式表明 \(U_\flat\) 在 \(S\) 中开.

\textup{\SGARefLink{concl:sga3-VIB-2.5-i}{(i)}} 的第二个断言由上述论证得到,
因为此时
\[
  Y=H,\qquad X=G,\qquad
  V_\flat=W_\flat=\pi_G^{-1}(U_\flat).
\]

证明断言 \textup{\SGARefLink{concl:sga3-VIB-2.5-ii}{(ii)}}.
设
\[
  V_{\mathrm{l.t.f.}}\supset V_{\mathrm{l.q.f.}}
\]
为 \(G\) 中使 \(u\) 为有限型 (分别地, 拟有限) 的点组成的开子集.
令
\[
  X=V_{\mathrm{l.t.f.}},
\]
并用 \(u_X\) 表示 \(u\) 在 \(X\) 上的限制.
由假设, \(X\) 包含 \(\epsilon(S)\).
设 \(x\in X\), 令 \(s=\pi_X(x)\).

若 \(u\) 在 \(x\) 处拟有限, 则由基变换可知 \(u_s\) 在 \(x\) 处拟有限;
又由于假设 \(G_s\) 局部有限型, 由
\SGARefLink{cor:sga3-VIB-1.3.1}{1.3.1}, \(u_s\) 局部拟有限.

反之, 若 \(u_s\) 局部拟有限, 则
\[
  u_X^{-1}\bigl(u_X(x)\bigr)
    \subset u_s^{-1}\bigl(u_s(x)\bigr)
\]
是有限集. 由于
\[
  u_X:X\longrightarrow H
\]
局部有限型, Chevalley 半连续性定理 (EGA~IV\(_3\), 13.1.3)
表明 \(x\) 在 \(X\) 中有一个开邻域 \(W\),
使对每个 \(x'\in W\), 纤维
\[
  u_X^{-1}\bigl(u_X(x')\bigr)
\]
离散. 因而由 EGA~II, 6.2.2
(以及 EGA~III\(_2\), Err\(_{\mathrm{III}}\)~20),
\(u_X\) 在 \(x\) 处拟有限.

因此,
\[
  U_{\mathrm{l.q.f.}}
    =\epsilon_G^{-1}\bigl(V_{\mathrm{l.q.f.}}\bigr),
  \qquad
  V_{\mathrm{l.q.f.}}
    =\pi_X^{-1}\bigl(U_{\mathrm{l.q.f.}}\bigr),
\]
第一个等式表明 \(U_{\mathrm{l.q.f.}}\) 在 \(S\) 中开.

若此外 \(u\) 在 \(S\) 上局部有限型,
则 \(X=G\), 因而 \(V_{\mathrm{l.q.f.}}\) 是
\(U_{\mathrm{l.q.f.}}\) 在 \(\pi_G\) 下的逆像.%
\begingroup
\renewcommand{\thefootnote}{\fnsymbol{footnote}}%
\footnote[1]{\label{note:sga3-VIB-2.5-u-loci}译者注:
在此句及下面关于 \(V_{\mathrm{n.r.}}\) 的类似句子中,
法文底本将 \(u\) 印作 \(G\).
陈述假设 \(u\) 局部有限型 (分别地, 局部有限呈示),
且 \(X\) 和 \(Z\) 正是 \(u\) 的相应性质成立的点集;
因此这里使用所需的 \(u\).}%
\endgroup{}
这证明了 \textup{\SGARefLink{concl:sga3-VIB-2.5-ii}{(ii)}} 的前半部分.

现考虑由使 \(u\) 为有限呈示 (分别地, 非分歧) 的点组成的开子集
\[
  V_{\mathrm{l.p.f.}}\supset V_{\mathrm{n.r.}},
\]
并假设
\[
  Z=V_{\mathrm{l.p.f.}}
\]
包含 \(G\) 的单位截面. 设 \(z\in Z\); 令
\[
  s=\pi_Z(z),\qquad h=u_Z(z).
\]

若 \(u\) 在 \(z\) 处非分歧, 则由基变换可知 \(u_s\)
在 \(z\) 处非分歧; 又由于假设 \(G_s\) 局部有限型,
由 \SGARefLink{cor:sga3-VIB-1.3.1}{1.3.1}, \(u_s\) 非分歧.

反之, 若 \(u_s\) 在点 \(z\) 处非分歧, 则纤维
\[
  u_s^{-1}(h)=u^{-1}(h)
\]
在点 \(z\) 处在 \(\kappa(h)\) 上非分歧.
由于 \(Z\) 是 \(G\) 的开子集, 纤维
\[
  u_Z^{-1}(h)=u_s^{-1}(h)\cap Z
\]
在点 \(z\) 处在 \(\kappa(h)\) 上非分歧.
由于 \(u_Z\) 局部有限呈示,
EGA~IV\(_4\), 17.4.1 蕴含 \(u_Z\) 在点 \(z\) 处非分歧.

因此
\[
  U_{\mathrm{n.r.}}
    =\epsilon_G^{-1}\bigl(V_{\mathrm{n.r.}}\bigr),
  \qquad
  V_{\mathrm{n.r.}}
    =\pi_Z^{-1}\bigl(U_{\mathrm{n.r.}}\bigr),
\]
第一个等式表明 \(U_{\mathrm{n.r.}}\) 在 \(S\) 中开.

若此外 \(u\) 在 \(S\) 上局部有限呈示,
则 \(Z=G\), 因而 \(V_{\mathrm{n.r.}}\) 是
\(U_{\mathrm{n.r.}}\) 在 \(\pi_G\) 下的逆像.
这完成了断言 \textup{\SGARefLink{concl:sga3-VIB-2.5-ii}{(ii)}} 的证明.

\medskip
\phantomsection\label{cor:sga3-VIB-2.6}%
\noindent\textbf{推论 2.6.---}
\emph{设 \(u:G\to H\) 是一个为根态射 (radicial) 的 \(S\)-群态射
(若 \(u\) 是单态射, 就属于这一情形; EGA~I, 3.5.4).
假设 \(G\) (分别地, \(H\)) 在 \(S\) 上局部有限呈示且平坦
(分别地, 局部有限型).
则使 \(u_s\) 为开浸入的 \(s\in S\) 组成的集合 \(U\)
在 \(S\) 中开, 且 \(u\) 在 \(U\) 上的限制为开浸入.}

由 \hyperref[concl:sga3-VIB-2.5-i]{2.5}\textup{\SGARefLink{concl:sga3-VIB-2.5-i}{(i)}},
使 \(u_s\) 平展的点 \(s\in S\) 组成的集合 \(U'\) 在 \(S\) 中开.
由于 \(u\) 是根态射, 对每个 \(s\in S\), \(u_s\) 也为根态射;
故由 EGA~IV\(_4\), 17.9.1, 有
\[
  U=U',
\]
这表明 \(U\) 开.
最后, 由 \hyperref[concl:sga3-VIB-2.5-i]{2.5}\textup{\SGARefLink{concl:sga3-VIB-2.5-i}{(i)}},
\(u\) 在 \(U\) 上的限制平展;
由于 \(u\) 是根态射, 此限制为开浸入 (EGA~IV\(_4\), 17.9.1).

\medskip
\phantomsection\label{prop:sga3-VIB-2.7}%
\noindent\textbf{命题 2.7.---}
\emph{设 \(G\) 是 \(S\)-群. 以下条件等价:}
\begin{enumerate}
\item[\textup{\SGAInlineRefTarget{prop:sga3-VIB-2.7:item:i}(i)}]
  \emph{\(G\) 在单位截面各点处在 \(S\) 上非分歧;}
\item[\textup{\SGAInlineRefTarget{prop:sga3-VIB-2.7:item:ii}(ii)}]
  \emph{单位截面为开浸入;}
\item[\textup{\SGAInlineRefTarget{prop:sga3-VIB-2.7:item:iii}(iii)}]
  \emph{\(G\) 在单位截面各点处在 \(S\) 上为有限呈示,
  且对每个 \(s\in S\), \(G_s\) 在 \(\kappa(s)\) 上非分歧.}
\end{enumerate}
\emph{若此外 \(G\) 在 \(S\) 上局部有限呈示,
则上述三个条件都等价于以下条件:}
\begin{enumerate}
\item[\textup{(iv)}] \emph{\(G\) 在 \(S\) 上非分歧.}
\end{enumerate}

断言 \textup{\SGARefLink{prop:sga3-VIB-2.7:item:i}{(i)}} 与
\textup{\SGARefLink{prop:sga3-VIB-2.7:item:ii}{(ii)}} 的等价性,
由下面更一般的 \hyperref[lem:sga3-VIB-2.7.1]{引理~2.7.1} 得到.
注意 (\SGARefLink{lem:sga3-VIB-1.3.1.1}{1.3.1.1}),
条件 \textup{\SGARefLink{prop:sga3-VIB-2.7:item:i}{(i)}} 与
条件 \textup{\SGARefLink{prop:sga3-VIB-2.7:item:iii}{(iii)}} 中的任何一个,
都蕴含对每个 \(s\in S\), \(G_s\) 在 \(\kappa(s)\) 上局部有限型.
因此, 由 EGA~IV\(_4\), 17.4.1,
断言 \textup{(i)} 等价于对每个 \(s\in S\),
\(G_s\) 在单位点 \(e_s\) (即 \(G_s\) 的单位元) 处在 \(\kappa(s)\) 上非分歧,
或由 \hyperref[cor:sga3-VIB-1.3.1]{1.3.1}, 等价于
\(G_s\) 在 \(\kappa(s)\) 上非分歧.
因此, 断言 \textup{\SGARefLink{cor:sga3-VIB-1.3.1:item:i}{(i)}} 与
\textup{\SGARefLink{prop:sga3-VIB-2.7:item:iii}{(iii)}} 等价.
最后, 若 \(G\) 在 \(S\) 上局部有限呈示,
则断言 \textup{(iii)} 与 \textup{(iv)} 等价 (EGA~IV\(_4\), 17.4.1).

\medskip
\phantomsection\label{lem:sga3-VIB-2.7.1}%
\noindent\textbf{引理 2.7.1.---}
\emph{设 \(G\) 是带有截面 \(\epsilon\) 的 \(S\)-概形.
\(G\) 在此截面各点处在 \(S\) 上非分歧的充分必要条件是
\(\epsilon\) 为开浸入.}

该条件的必要性来自 EGA~IV\(_4\), 17.4.1,
\textup{a)}\(\Rightarrow\)\textup{b\(^{\prime\prime}\)}).
反之, 若 \(\epsilon\) 为开浸入,
则结构态射
\[
  G\longrightarrow S
\]
在 \(\epsilon(S)\) 上的限制为同构;
故 \(G\) 在 \(\epsilon(S)\) 各点处在 \(S\) 上非分歧.
% Translation statistics: source_lines=321; translated_lines=307; inline_math=233; display_math=23; pending_TODO=0.
% Review: complete component translation; math fragments, reference keys, environments and note counts checked against source; formula prose translated; compilation pending parent integration.
